\section{Related work}\label{sec:related}

\paragraph{Long-video understanding.}
Benchmarks such as EgoSchema~\citep{mangalam2023egoschema},
Video-MME~\citep{fu2024videomme} and LongVideoBench~\citep{wu2024longvideobench}
show that accuracy of multimodal models degrades as videos grow longer and
that many questions require information spread over long temporal spans.
Architectural responses compress the visual token stream inside the model:
MovieChat~\citep{song2024moviechat} consolidates similar frame tokens into a
short- and long-term memory, and MA-LMM~\citep{he2024malmm} maintains a memory
bank of past visual features that is updated online. These memories are
\emph{implicit} (feature tensors) and part of a trained model; they cannot be
inspected record by record, and their contents carry no explicit timestamps.

\paragraph{Explicit memory for language agents.}
Retrieval-augmented generation~\citep{lewis2020rag} answers questions from an
external store queried at inference time. Generative
Agents~\citep{park2023generative} keep a natural-language memory stream scored
by recency, importance and relevance, and periodically write higher-level
reflections; MemGPT~\citep{packer2023memgpt} treats the context window as a
small main memory backed by external storage, with explicit operations for
moving information between the two. These systems motivate explicit,
inspectable records with salience-based eviction and summarisation, which this
project adopts, but they were designed for dialogue rather than for
perceptual streams in which the same physical object recurs and changes state.

\paragraph{Positioning.}
This project studies an \emph{explicit, structured, budgeted} memory for a
perceptual stream, with three properties not jointly present in the work
above: (i)~typed records for entities, events, state changes and relations;
(ii)~provenance from every record to source time intervals, preserved through
merging and compression, so that grounding can be measured; and (iii)~a
controlled comparison against simpler organisations under an identical budget
and identical observations. The emphasis is analytical: which information
survives, which mechanism is responsible for which gain, and where errors
originate.

\paragraph{Speech recognition.}
Transcripts are produced with Whisper~\citep{radford2023whisper} through the
CTranslate2-based faster-whisper implementation~\citep{fasterwhisper}. Whisper
is known to emit fluent but non-existent text on silence and non-speech
audio~\citep{koenecke2024careless}; the preprocessing choices
in \cref{sec:asr} are aimed at that failure mode.
